% 07_anexos/anexos.tex
% Anexos — soporte metodológico y documental de la tesis
% Las tablas extensas ya incorporadas en capítulos se referencian; no se duplican
% Instrumento de entrevista incluido como soporte metodológico; no como resultado empírico

\chapter{Anexos}

% ==============================================================
\section*{Anexo A. Cuadro de categorías, dimensiones, instrumentos y unidades de análisis (CDIU)}
\addcontentsline{toc}{section}{Anexo A. Cuadro de categorías, dimensiones, instrumentos y unidades de análisis (CDIU)}
% ==============================================================

El Cuadro de Categorías, Dimensiones, Instrumentos y Unidades de Análisis (CDIU) es el instrumento metodológico que organiza las categorías conceptuales centrales de la investigación, sus dimensiones jurídicas verificables, los instrumentos de análisis aplicados y las unidades de análisis examinadas en cada componente. Fue elaborado con base en las normas constitucionales, legales, jurisprudenciales y doctrinales identificadas durante la Fase~3 del proyecto y articula directamente los cuatro objetivos específicos de la investigación.

El cuadro identifica nueve categorías de análisis: derecho ambiental ecuatoriano, derechos de la naturaleza, actividad minera privada, control estatal ambiental, seguridad jurídica del sector minero privado, reformas legislativas y normativas mineras verificables, consulta ambiental y participación ciudadana, responsabilidad ambiental y reparación integral, y propuesta de compatibilidad constitucional y ambiental. Cada categoría se desagrega en dimensiones verificables que identifican el instrumento de análisis aplicado ---fichas documentales, matrices normativas, jurisprudencia vinculante, doctrina--- y la unidad de análisis correspondiente ---artículos constitucionales, disposiciones legales, holdings jurisprudenciales, principios del derecho ambiental ecuatoriano.

La versión sistematizada del CDIU con las nueve categorías y sus componentes esenciales se presenta en el Capítulo~I (Tabla~\ref{tab:cdiu}). La versión extendida con 35 filas y articulación detallada con los objetivos específicos fue desarrollada como instrumento interno del proceso de investigación y sirvió de base para la construcción de las matrices de análisis documental, normativo y jurisprudencial.

% ==============================================================
\section*{Anexo B. Matriz de análisis documental}
\addcontentsline{toc}{section}{Anexo B. Matriz de análisis documental}
% ==============================================================

La Matriz de análisis documental es el instrumento que organiza el corpus de fuentes verificadas de la investigación por tipo, identificación del documento, año, estado jurídico o documental y objetivo específico al que contribuye de manera primaria. Fue construida durante la Fase~3 de la investigación a partir de las 54 fichas documentales elaboradas para el repositorio del proyecto y contiene 24 entradas que representan las fuentes organizadas en ocho categorías: normativa, jurisprudencial, internacional, institucional, doctrinal, prensa contextual, reformas legislativas y un documento institucional contextual complementario.

La Matriz garantiza la trazabilidad del análisis: permite verificar que todas las dimensiones del análisis jurídico desarrollado en los Capítulos~II y~IV están respaldadas por fuentes verificadas incorporadas al repositorio del proyecto. Las columnas de la Matriz identifican el tipo de fuente, el instrumento o documento, el año de publicación, el estado jurídico o documental ---norma vigente, sentencia vinculante, tratado en vigor, informe institucional, doctrina académica, prensa contextual--- y el objetivo específico al que contribuye con mayor directitud.

La versión completa de esta Matriz se presenta en el Capítulo~III (Tabla~\ref{tab:analisis_documental}). Las restricciones documentales identificadas ---ausencia del expediente legislativo completo de la Ley de Fortalecimiento 2026 y de los proyectos de ley reformatorios tramitados entre 2018 y 2020--- se declaran expresamente en el Capítulo~III, Sección~3.7.3.

% ==============================================================
\section*{Anexo C. Matriz normativa del ordenamiento jurídico minero-ambiental ecuatoriano}
\addcontentsline{toc}{section}{Anexo C. Matriz normativa del ordenamiento jurídico minero-ambiental ecuatoriano}
% ==============================================================

La Matriz normativa sistematiza los diez instrumentos jurídicos que conforman el corpus normativo central de la investigación. Los instrumentos incorporados son:

\begin{enumerate}
    \item Constitución de la República del Ecuador (2008) — Registro Oficial No.~449 de 20 de octubre de 2008.
    \item Código Orgánico del Ambiente (2017) — Registro Oficial Suplemento No.~983 de 12 de abril de 2017.
    \item Reglamento al Código Orgánico del Ambiente (2019) — Registro Oficial Suplemento No.~507 de 12 de junio de 2019.
    \item Ley de Minería, versión codificada con reformas hasta agosto de 2025.
    \item Reglamento General a la Ley de Minería (2009).
    \item Ley Orgánica Reformatoria a la Ley de Minería (2013).
    \item Ley Orgánica para el Fortalecimiento de los Sectores Estratégicos de Minería y Energía (2026) — Registro Oficial Quinto Suplemento No.~234 de 2 de marzo de 2026. \textit{Advertencia: al momento de elaboración de esta tesis, esta ley acumula demandas de inconstitucionalidad ante la Corte Constitucional del Ecuador cuyo resultado no ha sido determinado.}
    \item Código Orgánico Integral Penal (2014) — Registro Oficial Suplemento No.~180 de 10 de febrero de 2014.
    \item Instructivo de Exploración y Explotación de Concesiones Mineras (2022).
    \item Resolución No.~0029-2025 de ARCOM sobre diferenciación entre minería formal, informal e ilegal.
\end{enumerate}

Cada instrumento está identificado con su institución emisora, Registro Oficial, artículos relevantes para la investigación y su nivel de relevancia para el análisis de la compatibilidad constitucional y ambiental de las reformas mineras. La Matriz permite verificar la trazabilidad de las afirmaciones jurídicas de los Capítulos~II y~IV con respecto a las normas citadas en el texto.

La versión sistematizada de esta Matriz se presenta en el Capítulo~II (Tabla~\ref{tab:normativa}).

% ==============================================================
\section*{Anexo D. Matriz jurisprudencial}
\addcontentsline{toc}{section}{Anexo D. Matriz jurisprudencial}
% ==============================================================

La Matriz jurisprudencial sistematiza los seis documentos del corpus jurisprudencial de la investigación, con especial atención a las tres sentencias vinculantes de la Corte Constitucional del Ecuador que establecen los parámetros de control constitucional aplicables al sector minero.

La \textbf{Sentencia No.~1149-19-JP/21} (noviembre de 2021) establece estándares vinculantes sobre los derechos de la naturaleza, el principio de precaución, el principio \textit{in dubio pro natura} y la consulta ambiental del artículo~398 de la Constitución, en el contexto del caso Bosque Protector Los Cedros. Sus holdings sobre la carga de la prueba de inocuidad, la nulidad e inejecutabilidad de las autorizaciones otorgadas sin consulta ambiental previa y la reparación integral de la naturaleza son determinantes para el análisis de compatibilidad del Capítulo~IV.

La \textbf{Sentencia No.~32-17-IN/21} (diciembre de 2021) establece el principio de reserva de ley orgánica como límite constitucional a las normas reglamentarias que pretendan regular materias con incidencia en derechos constitucionales en el sector minero. Su holding sobre los límites de la potestad reglamentaria es aplicable directamente al análisis del criterio 7 (consulta previa, libre e informada) de la Matriz del Capítulo~IV.

La \textbf{Sentencia No.~22-18-IN/21} (2021) distingue con precisión la consulta ambiental del artículo~398 de la Constitución de la consulta previa, libre e informada (CPLI) del artículo~57.7, establece el contenido mínimo de la consulta ambiental con referencia al Acuerdo de Escazú, y declara inconstitucionales los artículos~462 y~463 del Reglamento al Código Orgánico del Ambiente por regular la CPLI a nivel reglamentario en violación de la reserva de ley orgánica. El vacío normativo generado por esta declaratoria de inconstitucionalidad permanece sin subsanarse al momento de elaboración de esta tesis.

La Matriz incluye además la Guía de Jurisprudencia Constitucional sobre Derechos de la Naturaleza del CEDEC (2023) y el Comunicado institucional de la Corte Constitucional del Ecuador de agosto de 2025 sobre admisión de demandas y suspensión provisional de normas, incorporados como fuentes contextuales y de sistematización.

La versión sistematizada de esta Matriz se presenta en el Capítulo~II (Tabla~\ref{tab:jurisprudencial}).

% ==============================================================
\section*{Anexo E. Matriz de Verificación de Compatibilidad Constitucional y Ambiental para Proyectos de Ley de Apertura Minera Privada}
\addcontentsline{toc}{section}{Anexo E. Matriz de Verificación de Compatibilidad Constitucional y Ambiental para Proyectos de Ley de Apertura Minera Privada}
% ==============================================================

La Matriz de Verificación de Compatibilidad Constitucional y Ambiental para Proyectos de Ley de Apertura Minera Privada es la propuesta central del Capítulo~IV de la investigación. Consta de dieciséis criterios de evaluación jurídica organizados en cuatro grupos temáticos, cada uno con fundamento constitucional, legal o jurisprudencial verificable, estándar exigible mínimo e indicador de verificación.

\textbf{Grupo I. Criterios constitucionales de base} (criterios 1--5): derechos de la naturaleza como sujeto jurídico (arts.~71--74 CRE); principio de prevención (art.~396 CRE, art.~9 COA); principio de precaución (art.~73 CRE, Sentencia 1149-19-JP/21); principio de no regresividad ambiental (art.~11.8 CRE, art.~161 COA); principio \textit{in dubio pro natura} (art.~395.4 CRE).

\textbf{Grupo II. Criterios de participación y consulta} (criterios 6--8): consulta ambiental (art.~398 CRE, Sentencia 1149-19-JP/21); consulta previa, libre e informada (art.~57.7 CRE, Sentencia 22-18-IN/21); participación ciudadana ambiental (Acuerdo de Escazú, art.~18 CRE).

\textbf{Grupo III. Criterios de control y responsabilidad} (criterios 9--13): licenciamiento ambiental (arts.~162, 166 COA); fiscalización estatal efectiva (OC-23/17 CorteIDH); transparencia activa (Acuerdo de Escazú); responsabilidad ambiental objetiva del concesionario (arts.~396, 397 CRE); reparación integral (art.~397 CRE, Sentencia 1149-19-JP/21).

\textbf{Grupo IV. Criterios de compatibilidad sistémica} (criterios 14--16): seguridad jurídica del sector minero formal (art.~82 CRE); diferenciación entre minería privada formal, informal e ilegal (arts.~260, 261 COIP); incidencia constitucional en el sector minero privado (art.~316 CRE).

La aplicación ilustrativa de la Matriz a la Ley Orgánica para el Fortalecimiento de los Sectores Estratégicos de Minería y Energía de 2026, con indicación del nivel de cumplimiento de cada criterio, se desarrolla en el Capítulo~IV, Sección~4.2.2. Este análisis tiene carácter académico y no prejuzga el pronunciamiento de la Corte Constitucional del Ecuador en los procesos de control constitucional en curso.

La versión sistematizada completa de la Matriz, con las cinco columnas de evaluación, se presenta en el Capítulo~IV (Tabla~\ref{tab:compatibilidad}).

% ==============================================================
\section*{Anexo F. Instrumento de entrevista semiestructurada (previsto, no aplicado)}
\addcontentsline{toc}{section}{Anexo F. Instrumento de entrevista semiestructurada (previsto, no aplicado)}
% ==============================================================

\textbf{Nota metodológica.} Las entrevistas semiestructuradas previstas en el diseño de la investigación no fueron aplicadas durante el período de elaboración del trabajo de titulación, por restricciones de viabilidad operativa. El instrumento que se presenta a continuación fue diseñado como parte del diseño metodológico de la investigación y se incluye en esta sección como documento de soporte, no como resultado empírico. Su reproducción responde al requisito de transparencia metodológica: permite verificar que el instrumento estuvo efectivamente planificado y que su no aplicación no obedece a una omisión de diseño. En caso de que la investigación sea continuada o replicada, el instrumento puede ser aplicado en su estado actual sin modificaciones al diseño metodológico general.

\bigskip

\noindent\textbf{Instrumento No.~1. Guía de entrevista semiestructurada}

\medskip

\noindent\textit{Tema:} Compatibilidad de las reformas de apertura del sector minero privado con el derecho ambiental ecuatoriano.

\smallskip

\noindent\textit{Destinatarios:} Abogados especialistas en derecho ambiental o derecho minero; docentes universitarios de derecho ambiental o constitucional; funcionarios del Ministerio del Ambiente, Agua y Transición Ecológica (MAATE); funcionarios de la Agencia de Regulación y Control Minero (ARCOM); actores vinculados al sector minero privado formal.

\smallskip

\noindent\textit{Modalidad prevista:} Entrevista individual, presencial o virtual, con duración aproximada de cuarenta y cinco a sesenta minutos.

\smallskip

\noindent\textit{Instrucciones:} Las preguntas son orientativas. El orden puede modificarse según el curso de la conversación. El entrevistador puede formular preguntas de seguimiento ante respuestas que requieran mayor desarrollo.

\bigskip

\noindent\textbf{Preguntas:}

\begin{enumerate}

    \item ¿Cómo evalúa la compatibilidad de la Ley Orgánica para el Fortalecimiento de los Sectores Estratégicos de Minería y Energía, publicada en el Registro Oficial de 2 de marzo de 2026, con los estándares constitucionales de protección ambiental y derechos de la naturaleza establecidos en los artículos 71 al 74 y 395 al 399 de la Constitución de la República del Ecuador?

    \item ¿Considera que los cambios al régimen de autorización ambiental para actividades mineras introducidos por la Ley de Fortalecimiento 2026 fortalecen o debilitan el control estatal sobre las actividades mineras privadas, y qué efectos jurídicos concretos pueden derivarse de esa modificación?

    \item ¿Qué efectos jurídicos concretos genera el vacío normativo de la consulta previa, libre e informada en el sector minero ---derivado de la Sentencia No.~22-18-IN/21 de la Corte Constitucional del Ecuador--- sobre las concesiones mineras activas y futuras en territorios donde habitan comunidades, pueblos y nacionalidades indígenas?

    \item ¿En qué medida la apertura del catastro minero y la simplificación de procedimientos administrativos en las concesiones mineras resultan compatibles con el principio de no regresividad ambiental establecido en el artículo 161 del Código Orgánico del Ambiente, y qué criterios jurídicos permitirían determinar cuándo una simplificación procedimental constituye regresión ambiental?

    \item ¿Qué mecanismos institucionales considera indispensables para garantizar que la fiscalización estatal del sector minero privado sea efectiva, habida cuenta de las brechas en la capacidad de fiscalización disponible reportadas en el Informe de Rendición de Cuentas de ARCOM para el año 2024?

    \item ¿Qué lineamientos jurídicos permitirían armonizar la seguridad jurídica del inversionista minero privado ---entendida como previsibilidad normativa y estabilidad del régimen de concesiones--- con los derechos de la naturaleza, los estándares constitucionales de consulta ambiental y la obligación de reparación integral por daños ambientales en el Ecuador?

\end{enumerate}

